Se estudiaron dos fenómenos relacionados con el transporte de electrones en gases. En la primera parte se investigó el efecto Ramsauer--Townsend utilizando un tiratrón de xenón 2D21, registrando los voltajes de ánodo y de rejilla-blindaje como función del voltaje de cátodo para distintos voltajes de filamento. Las mediciones se realizaron tanto de manera manual como mediante un sistema automatizado, permitiendo comparar ambos métodos de adquisición de datos. En la segunda parte se empleó la técnica de Townsend pulsada para determinar la velocidad de deriva electrónica en una mezcla de 99\,\% Ar y 1\,\% N$_2$. A partir del tiempo de tránsito de un enjambre electrónico se calculó la velocidad de deriva para distintos valores del campo eléctrico reducido $E/N$. Los resultados obtenidos reproducen el comportamiento esperado para el transporte electrónico en gases y muestran la relación entre los procesos microscópicos de dispersión y las propiedades macroscópicas de transporte.